\documentclass[11pt]{article}
\usepackage[margin=1in]{geometry}
\usepackage{amsmath,amssymb}
\usepackage{enumitem}
\usepackage{hyperref}

\title{2023 F=ma Exam Review}
\author{Nathan P.}
\date{February 1, 2026}

\begin{document}
\maketitle

\section*{1}
\begin{itemize}[leftmargin=*, itemsep=0.25em]
\item Average velocity $=\Delta x/t$. Initial position was $2.25$ rotations ago. Two rotations ago it was in the same spot at the bottom; another fourth rotation ago it would be at left side of circle, so average velocity would be leftward and downward.
\end{itemize}

\section*{4}
\begin{itemize}[leftmargin=*, itemsep=0.25em]
\item The work the person does $=mgh+W_f$, and $W_f=\mu mg\cos\theta\cdot \dfrac{h}{\sin\theta}$.
\end{itemize}

\section*{6}
\begin{itemize}[leftmargin=*, itemsep=0.25em]
\item Look at the answer options---the height is $\ll$ hypotenuse, so use small angle approximations and relate the triangles of tension's components and the physical lengths triangle $(l,r,h)$.
\end{itemize}

\section*{8}
\begin{itemize}[leftmargin=*, itemsep=0.25em]
\item ``Rest frame'' means $v_{cm}=0$, so total momentum $=0$, so $m_a v_a=m_b v_b$ so $v_a/v_b=m_b/m_a=1/10$.
\end{itemize}

\section*{10}
\begin{itemize}[leftmargin=*, itemsep=0.25em]
\item Use the range formula $R=\dfrac{v^2\sin 2\theta}{g}$.
\end{itemize}

\section*{11}
\begin{itemize}[leftmargin=*, itemsep=0.25em]
\item Combine range and height formulas to get that $H=R$ when $\tan\theta=4$.
\item Max height happens when $\theta=90$ and $H=v^2/2g$ which $=h/\sin^2(\arctan 4)$.
\end{itemize}

\section*{13}
\begin{itemize}[leftmargin=*, itemsep=0.25em]
\item This one tests your physics sense. For the box to lift, it will reach $N=0$, which means somewhere along the way friction $(\mu N)$ will be not enough to stop sliding, so it will always slide first (before $N=0$).
\end{itemize}

\section*{15}
\begin{itemize}[leftmargin=*, itemsep=0.25em]
\item Good job recalling $v_{\text{orb}}=\sqrt{GM/r}$ and since $v$ appears constant, $M\propto r$.
\end{itemize}

\section*{16}
\begin{itemize}[leftmargin=*, itemsep=0.25em]
\item Since it is a small angle pendulum, use SHM equation $\theta(t)=\theta\cos\omega t$.
\item Solve for $t$ when $\theta(t)=\theta\cos\omega t=-\theta/2$.
\item Make sure to recognize that is half the time between bounces.
\end{itemize}

\section*{20}
\begin{itemize}[leftmargin=*, itemsep=0.25em]
\item It describes the picture for your imagination. Mass attached to end of string to form a small angle pendulum. It has $v_x$ and $v_y$.
\item $v_x=0$ and $v_y=0$ at top of trajectory so $(0,0)$ should be on graph.
\item Count how many times it should hit zero: two times per period for $v_x$, four times for $v_y$, so they should cross the $y$ and $x$ axis $2$ and $4$ times respectively.
\end{itemize}

\section*{21}
\begin{itemize}[leftmargin=*, itemsep=0.25em]
\item Need clarification on this problem but from what I'm getting, $\sum F=0$ so $v_{cm}=\text{constant}$ and $\Delta p=0$ and both ring and mass have mass of $m$ so they are moving at $v/2$ relative to each other, then angular momentum conservation confirms $r$ is at least $R/2$ but geometry limits it to max $R/2$ (stays inside ring) so it does travel in circle of radius $R/2$ about the cm. $T=2\pi(R/2)/(v/2)=2\pi R/v$.
\end{itemize}

\section*{22}
\begin{itemize}[leftmargin=*, itemsep=0.25em]
\item They are different spring constant $k$'s. Remember for springs in series: $\frac{1}{k_{\text{eff}}}=\frac{1}{k_1}+\frac{1}{k_2}+\cdots$. $k_B$ is tricky but it doubles that force value because pulling $x$ extends spring by $2x$ and it doubles that value again because the forces add on both sides (and they are the same force).

\vspace{0.5em}
\begin{center}
$AM\ge GM\ge HM\quad \dfrac{a+b}{2}>\sqrt{ab}>\dfrac{2ab}{a+b}$
\end{center}
\end{itemize}

\section*{24}
\begin{itemize}[leftmargin=*, itemsep=0.25em]
\item Essentially, is it more likely for the ball to hit the plate when it is moving up or down? The answer is up because let's say the ball reached the collision zone's border. If the plate was moving up at any location, the ball will collide while plate moves up, so it will increase its speed. However, if plate is moving down and is at a low enough region, it has time to turn around and collide with the ball while moving up.
\end{itemize}



\end{document}